\subsection{Quantifiers, Bound Variables, and Negation}

\begin{principle} \label{princ:analysis-language-as-structure} \hlt{Language Carries Structure}\\
Objects, hypotheses, conclusions, and quantifier order matter. For instance,
\begin{align}
\forall x\in X\ \exists y\in Y\ P(x,y), \qquad
\exists y\in Y\ \forall x\in X\ P(x,y) \nonumber
\end{align}
are different: in the first $y$ may depend on $x$; in the second one $y$ works for all $x$.
\end{principle}

\begin{method} \label{meth:proof-grammar-and-negation} \hlt{Proof Grammar and Negation}\\
In theorem statements, put the hypotheses before the conclusion and keep quantifiers in the order meant.\\
In particular, switching $\forall$ and $\exists$ usually changes the statement; see Principle \ref{princ:analysis-language-as-structure}.\\
Before applying a rule, also check that the displayed expressions are well formed; an undefined expression makes the sentence ill formed, not merely false.\\
The word ``or'' is inclusive in mathematics, and the basic negation rules are
\begin{align}
\neg(a\land b)\Longleftrightarrow \neg a\lor\neg b, \qquad \neg(a\lor b)\Longleftrightarrow \neg a\land\neg b, \qquad \neg(a\Rightarrow b)&\Longleftrightarrow a\land\neg b \nonumber
\end{align}
Thus a proof by contradiction of $a\Rightarrow b$ assumes both $a$ and $\neg b$ and derives an impossibility.\\
\hlt{Key Insight}: Before proving or negating a statement, parse its quantifiers and logical connectives.
\end{method}

\hlt{Placeholder}: Add negation rules for quantified statements and examples translating English into symbolic logic.
